\section{Cover-Mediated Feedback Burden}

\subsection{Feedback burden from cover-mediated readout}

The base bridge count
\[
B=137
\]
is not the final effective bridge capacity. The compact fiber also
supplies cover-mediated feedback from the paired magnetic sector into
the electromagnetic bridge.

The paired magnetic sector is
\[
Z_4\times Z_4,
\]
and the electric-cover sector is
\[
Z_8.
\]
Since
\[
N_m\mid N_e,
\qquad
8=2\cdot4,
\]
the compact fiber contains the canonical electric-cover / magnetic-base
quotient
\[
\pi:Z_8\to Z_4.
\]
This quotient is derived from the sector-capacity relation and is not
independent bridge data.

The target observable
\[
\alpha_F
\]
is a carrier-level electromagnetic bridge coupling. Paired magnetic
feedback therefore does not enter as a purely magnetic-internal quantity;
it is read through the electric-cover / magnetic-base interface.

\subsection{Canonical cover-interface readout routes}

The base-cover feedback channel has two mechanism-distinct scalar
readout resolutions. The first is direct electric-cover resolution,
\[
\frac{1}{N_e}.
\]
The second is fully resolved joint magnetic/electric resolution,
\[
\frac{1}{N_mN_e}.
\]
These are not sequential filters applied to the same event. Sequential
transport would multiply capacities; mechanism-distinct scalar access
contributions to a single feedback burden add before the burden is
distributed over the final bridge capacity.

Thus the gross cover-mediated readout is
\[
\frac1{N_e}+\frac1{N_mN_e}.
\]
At common denominator \(N_mN_e\),
\[
\frac1{N_e}+\frac1{N_mN_e}
=
\frac{N_m}{N_mN_e}
+
\frac1{N_mN_e}
=
\frac{N_m+1}{N_mN_e}.
\]
The numerator
\[
N_m+1
\]
is the intrinsic mechanism count of the base-cover feedback channel.
For
\[
N_m=4,
\]
this gives
\[
N_m+1=5.
\]
Therefore the gross feedback burden is
\[
\boxed{
M_{\rm fb}=5.
}
\]

The local product-fiber denominator
\[
N_mN_e
\]
is appropriate when the target observable is resolved at product-fiber
scale. That is not the normalization used for the alpha bridge equation.
The target observable is the effective bridge capacity
\[
x=\alpha_F^{-1},
\]
so the bridge share is
\[
\frac1x.
\]
Reusing the local denominator \(N_mN_e\) and then dividing again by
\(x\) would double-count target normalization. For the alpha bridge, the
feedback contribution has the form
\[
\frac{K}{x},
\]
not
\[
\frac{K}{N_mN_e\,x}.
\]
Thus the intrinsic feedback numerator \(N_m+1\) is carried forward as
the gross feedback burden.

\subsection{Scalar self-resolution cost}

The paired magnetic contrast space is
\[
G=Z_4\times Z_4,
\]
with
\[
|G|=N_m^2=16.
\]
The datum state
\[
(0,0)
\]
carries no paired magnetic contrast and is excluded from the active
contrast support. Define
\[
G^\ast=G\setminus\{(0,0)\}.
\]
Then
\[
|G^\ast|=N_m^2-1.
\]
For \(N_m=4\),
\[
|G^\ast|=16-1=15.
\]

The electromagnetic bridge coupling is a scalar carrier-level observable.
It does not resolve the individual labels of \(G^\ast\). The unresolved
non-datum contrast support therefore contributes one scalar
self-resolution cost distributed over the active contrast space:
\[
A_{\rm self}
=
\frac1{|G^\ast|}
=
\frac1{N_m^2-1}.
\]
For \(N_m=4\),
\[
A_{\rm self}
=
\frac1{15}.
\]
This is a self-resolution cost, not an additional feedback source and
not an attenuation factor applied independently to every feedback event.

\subsection{Closure-budget accounting}

The closure budget contains the gross generated feedback burden
\[
M_{\rm fb}=N_m+1
\]
and the consumed self-resolution cost
\[
A_{\rm self}=\frac1{N_m^2-1}.
\]
The source-level budget relation is
\[
\text{gross generated feedback}
=
\text{delivered feedback}
+
\text{consumed self-resolution cost}.
\]
Equivalently,
\[
\text{delivered feedback}
=
\text{gross generated feedback}
-
\text{consumed self-resolution cost}.
\]
Thus
\[
K=M_{\rm fb}-A_{\rm self}.
\]

This excludes a multiplicative form such as
\[
M_{\rm fb}(1-A_{\rm self})
\]
in the present role assignment. A multiplicative form would treat
\(A_{\rm self}\) as an attenuation applied to each feedback event. Here
\(A_{\rm self}\) is an additive cost consumed from the same closure
budget before the net burden is delivered to the bridge.

Substituting the compact-fiber values,
\[
K
=
(N_m+1)
-
\frac1{N_m^2-1}.
\]
For \(N_m=4\),
\[
K
=
5-\frac1{15}
=
\frac{74}{15}.
\]
Therefore the net first-order feedback burden is
\[
\boxed{
K=\frac{74}{15}.
}
\]

\subsection{\texorpdfstring{Role of \(K\) in the bridge equation}{Role of K in the bridge equation}}

The quantity
\[
K=\frac{74}{15}
\]
is not a direct channel count to be added to \(B\). It is a feedback
burden delivered into the effective bridge capacity. Therefore it appears
as
\[
\frac{K}{x}
\]
in the self-consistent bridge equation, not as
\[
K.
\]
The first-order equation is
\[
x=B+\frac{K}{x},
\]
not
\[
x=B+K,
\]
and not merely a one-pass correction
\[
x=B+\frac{K}{B}.
\]
The denominator is the final effective bridge capacity \(x\), because
the feedback modifies the same bridge capacity over which it is
distributed.

The feedback-burden dependency structure is therefore:
\[
\text{cover-mediated gross burden}
\quad
N_m+1=5,
\]
\[
\text{active non-datum contrast support}
\quad
N_m^2-1=15,
\]
\[
\text{scalar self-resolution cost}
\quad
\frac1{N_m^2-1}=\frac1{15},
\]
\[
\text{net delivered feedback}
\quad
K=(N_m+1)-\frac1{N_m^2-1}.
\]
Thus
\[
\boxed{
K=\frac{74}{15}
}
\]
is a structural feedback burden determined by the compact-fiber closure
budget, not a fitted constant.